\section{Supplying the common comparison estimates}
\label{sec:parameter-comparison-suppliers}
The preceding comparison arguments use one collection of local
bounds for the base representation. We show how the local estimates
already proved provide that collection, and how its parameters are
chosen before a running representation is inserted. Constants in a
fixed-volume continuity estimate may depend on that volume. The
radius of the allowed coarse variations and the large-coupling
threshold used for the construction have a stronger uniformity.

Fix the branch data, a positive target support reserve, the derivative
order $j$, the cutoff convention, and positive width and packet-energy
parameters. The scale factor is taken above the common floor in the
local activity and Gaussian estimates. Fix a coarse variation scale
$v(x)\ge0$ of upper exponent $-1/2$. All estimates in this section
refer to the same fine representation, sector, tags and held-field
compensation.

\subsection{Elementary factors and conditional packet actions}
The construction in Section~\ref{sec:activity-jet-assembly} provides
prices for the elementary compensated factors together with their
support-dependent derivative scales. At the fixed derivative order,
choose a positive spare support and extra-cost reserve $\delta$.
If $s$ denotes the support-cover plus length charge, then
$(1+s)^j\le C_{j,\delta}e^{\delta s}$ for $s\ge0$: the continuous
function $(1+s)^je^{-\delta s}$ is bounded on $[0,\infty)$.
Apply the weighted price estimate with this reserve before removing
the weights. This pays the factor $r(x)^j(1+s)^j$ from the derivative
scale and gives scale-one jet prices at every order through $j$.
The warning catalogue is finite on the fixed torus and its finite
jet prices can also be summed. To pass
from a rooted weighted sum to its unweighted sum on a fixed torus,
sum the rooted inequality over all its sites. Every nonzero factor
has nonempty assigned support and is counted at least once; all
removed support and extra-cost weights are at least one. The same
argument applies to the base logarithmic prices. It gives both the
finite total elementary price and the finite common bound on the
absolute values of the real logarithms used in
Section~\ref{sec:local-parameter-comparison}. The prices need not be
small after this volume-dependent summation.

There is a separate condition on the packet action: its differentiated
bound is needed whenever the enforcement cutoff is not a zero germ.
Write that cutoff as a finite product of its bond cutoffs. If any
bond cutoff were zero on a neighborhood of the base point, the whole
product would be zero there. Thus a nonzero packet germ makes every
relevant bond cutoff a nonzero germ. The compensated cutoff estimates
then bound the corresponding increments and their derivatives by
\[
                    s(x)+2a\tau(x),
\]
where $s$ has upper exponent $-1/2$ and $a>0$ is the fixed
compensation radius parameter. The mean coordinates have the same
upper exponent, by Section~\ref{sec:mean-action-jets}. Enlarge these
bounds to one nonnegative function of upper exponent zero.

The last subsection of Section~\ref{sec:mean-action-jets} applies to
an arbitrary held fine field. Each local packet-action summand is
a held fine action, a varying coarse action, or an extracted affine
square. Its derivatives are bounded by the corresponding local
constants there. For a packet with $s$ blocks, the number of such
energy addresses is at most
\[
                  s\,b(200B+1)\,n_L,
\]
where $b(r)$ is the finite lattice-ball count at radius $r$ and $n_L$
is the fixed number of fine energy addresses per coarse block.
Summing the local jet bounds, including the far-field term, gives
$A(x)s$ with one finite $A(x)\ge0$. The same bound holds on every
admissible chart and for every held fine field. In particular this
supplies the conditional packet-action estimate, without requiring
smoothness of the fine input or an upper bound on its uncut energy.

\subsection{Every read subset has the same polynomial bound}
Let $I$ be any subset of fine bonds, including the empty subset.
The projected covariance contains all colours of those bonds.
Write $C_I(y)$ for this covariance and $w_I(y,z)$ for the compensated
Gaussian coordinate. Its inverse exists by the common positive
covariance estimate. The quadratic exponent is
\[
              q_I(y,z)=\tfrac12 w_I(y,z)^T C_I(y)^{-1}w_I(y,z).
\]
The inverse-matrix estimates and weighted matrix product estimates
of Sections~\ref{sec:normalization-jets} and
\ref{sec:gaussian-differentiation} apply to every principal read
subset with the same row constants. This is where the projected
covariance bounds, rather than merely positivity of the full
covariance, are used.

For clarity, the finite-dimensional multiplication giving the
quadratic bound can be seen directly. Suppose the coordinate jets
of $w_I$ are at most $W$, and the jets of the inverse matrix have
row sums at most $B_r$ at order $r\le j$. A derivative of
$w_I^TC_I^{-1}w_I$ distributes its ordered directions among three
factors, giving at most $3^j$ terms. Sum the middle matrix row first.
Each term is at most $|I|\dim\mathfrak g$ times
$W^2\max_{r\le j}B_r$. The factor $1/2$ can be discarded in an upper
bound. On the full torus the compensated-coordinate estimates give
\[
 W\le \left[1+g(\beta)^{-1}
           \{s(x)+2A_{\rm mean}(x)\}\right]F_0(z)=V_0F_0(z),
\]
where $F_0$ uses every fine bond and colour and $V_0\ge1$ is finite.
Since $|I|\dim\mathfrak g$ is bounded by the full Gaussian dimension,
we obtain
\begin{equation}\label{eq:comparison-quadratic-supplier}
                 \|q_I\|_{j,y_0}\le C_*F_0(z)^2
                         \quad\hbox{for every }I.
\end{equation}
All constants on the right are independent of the chosen subset.
For $I=\varnothing$, the quadratic form is zero.

The relative density compensation $K_I$ is the Jacobian from
Section~\ref{sec:compensated-variation} times the Gaussian density
ratio from Section~\ref{sec:gaussian-differentiation}. At the base
point its value is one. The finite-order determinant and inverse
estimates there give, for some integer $k$ and nonnegative finite
$U(x)$, a derivative scale bounded by
$U(x)F_I(z)^k$, where
\[
              F_I(z)=1+|I|\dim\mathfrak g
                          +\sum_{(e,a)\in I\times\mathfrak g}|z_{e,a}|.
\]
In this expression the colour index means a fixed orthonormal basis.
The bound uses only the prescribed derivative order, the local row
prices and the compensation Jacobian bounds. Its polynomial degree
is fixed before the subset is selected. Since $1\le F_I\le F_0$,
\[
 \max_{0\le r\le j}\{U(x)F_I(z)^k\}^r
        \le(1+U(x))^jF_0(z)^{kj}.
\]
Together with \eqref{eq:comparison-quadratic-supplier}, this supplies
\eqref{eq:local-parameter-density-input} simultaneously for all $I$,
with $N=kj$ and
$C_0=(1+U(x))^j+C_*$. The quadratic constant $C_*$ is allowed to
depend on the torus and base coupling; no uniform volume bound is
asserted for it.

The same compensated map provides the held-field and protected-patch
identities. The open-neighborhood envelope in
Section~\ref{sec:gaussian-differentiation} justifies differentiation
of its integrals. Thus the algebraic pointwise comparison uses the
actual integrands to which that differentiation theorem applies.

\subsection{Unrestricted packet values}
The lower packet-energy bound in the seed estimates applies on the
support of the packet mask. Adding the nonnegative far-field energy
preserves this lower bound. At a fixed stage it gives
$u\ge-E(x)s\ge-E(x)s_*$ in the notation of
\eqref{eq:local-parameter-forced-packet}. The product of profile,
charge and enforcement factors lies in $[0,1]$. Its remaining ticket
factor is bounded by $\exp(|c_{\rm ticket}|s_*)$, so this number can
be taken as $M_0$.

Each condensation value has a finite base bound, and each local
Haar-chart correction has a finite absolute bound. There are at most
as many terms as fine bonds. The remaining normalization term is
$-n_s\log m(\beta)$, again with $n_s$ no larger than that finite bond
count. Their sum gives a finite $A_0$ bounding $v_\beta$ from above.
Only this last normalization term changes between the compared
representations; hence
$|v_{\beta'}-v_\beta|\le n_*t$. These are exactly the hypotheses of
the damping estimate already proved. Positive unbounded energies
are controlled by $\beta'\ge\beta/2$ and require no further cutoff.

\subsection{Choosing the common variation scale and thresholds}
We spell out the quantifier order needed for a running history.
The common scale-factor floor is chosen from the elementary-activity
and Gaussian local estimates before the derivative order, cutoff
convention and positive width parameters are supplied. For a scale
factor above that floor and fixed branch data, the compensation
chart has a fixed positive radius. Its protected-field requirement
is eventually satisfied because $\tau(x)\to0$.

The elementary activity estimate is applied at the prescribed scale
$v(x)$ and the compensated-increment scale $s(x)$, both of upper
exponent $-1/2$. The packet-action, Jacobian, density and quadratic
estimates just established have eventual validity at those same
scales. Take the intersection of their finitely many eventual
conditions. This fixes one threshold in $x$ before the finite tori
and the incoming representation are supplied. The associated
constants and summable prices are then instantiated at that actual
representation. In particular, a later fixed-volume comparison
neighborhood may depend on its total prices and covariance, as the
continuity proof permits.

To combine this with the common integrated derivative estimate,
choose a fixed $Z>0$ large enough for the normalized exponential
coarse variations at derivative order $j$. Set
\[
                    t_0(x)=2\tau(x),\qquad
                    v(x)=(j+1)Zt_0(x).
\]
This function is eventually positive and has upper exponent $-1/2$.
The elementary inequality $jZt_0\le v$ pays the ordered variation
directions. Apply both the local comparison construction and the
integrated Gaussian estimate at this same $v$. The latter supplies
the ordinary component price and derivative-scale data. The cluster
budget construction, with the strictly positive target support
reserve, then supplies the ordinary cluster criterion with its
strict spare margin. All are evaluated at the same incoming data;
no alternate representation is substituted between these steps.

Finally, a schedule has only finitely many transitions at each fixed
stage. Take the maximum of their scale-factor floors at that stage.
After those factors are fixed, take a common upper threshold for the
finitely many transition conditions, enlarge the stage entry coupling
to that threshold, and preserve any earlier lower bound. Increasing
entry couplings further preserves these conclusions. Consequently the comparison data and cluster margin are available
uniformly for each already supplied transition satisfying the
underlying sector, shell-access, protected-chart and incoming activity
hypotheses. This selection does not construct those incoming data.

For an application to a running curve one must also retain its
renewed price bounds, the geometric bounds on the downward map,
the admissible initial action, and the inequality
$\operatorname{preExponent}/B\le\operatorname{width}/2$.
The curve has the same clipping and path data and the same regular
and forced dictionaries. Its displayed output must equal the
constructed nonlinear output under the chosen dictionary
identifications. With these hypotheses, the fixed-volume theorem
applies to that actual output, giving regular jet and forced value
control and its subsequential summable majorants. The construction
of an entire history satisfying all these hypotheses remains a
separate step.
